\documentclass[12pt,letterpaper]{article}

\usepackage[utf8]{inputenc}
\usepackage{natbib}
\usepackage{graphicx}
\usepackage{indentfirst}
\usepackage{amsmath}
\usepackage{amsfonts}
\usepackage{amssymb}
\usepackage{comment}
\usepackage{url}
\usepackage[left=4cm,right=2.5cm,top=4cm,bottom=2.5cm]{geometry}

\setlength{\parindent}{2cm}
\renewcommand{\baselinestretch}{1.5}
\renewcommand\contentsname{Índice de contenidos}
\renewcommand\refname{Referencias}
\renewcommand\figurename{Figura}

\usepackage{graphicx}
\begin{document}

\newpage
\vspace*{-.5cm}
\begin{picture}(18,4)(0,30)
	\put(350,-20){\includegraphics[scale=0.25]{./images/LogoUsach.pdf}}
\end{picture}

\sloppy
\thispagestyle{empty}
\vspace*{-1.6cm}

\begin{center}
	{\bf \mbox{\large UNIVERSIDAD DE SANTIAGO DE CHILE}}\\
	{\bf \mbox{FACULTAD DE INGENIER\'IA}}\\
	{\bf \mbox{DEPARTAMENTO DE INGENIER\'IA INFORM\'ATICA}}\\
\end{center}

	\vspace*{2cm}
	\par
	\begin{center}
	\large
		\textsl{Paradigmas de Programaci\'on}
	\end{center}

	\vspace{0.5cm}
	\begin{center}
	\large
		\textbf{Informe N$^{\circ}$2\\
		Sistema de gestión fastClinic}
	\end{center}
	
	\vspace{3cm}
	\begin{flushright}
		\begin{tabular}[t]{l l}
			Nombre: & Joaquín Ignacio Villagra Pacheco \\
			R.U.N.: & 18.847.934-0\\
			Carrera: & Ingeniería Civil Informática\\
			A\~no estimado de egreso: & 2018\\
			E-mail: & joaquin.Villagra@usach.cl\\
			Profesor: & José Allende\\
			Ayudante: & Maurcio Rojas

		\end{tabular}
	\end{flushright}
	\begin{center}
		\vspace{1.5cm}
		Santiago, Chile\\
		Viernes, 17 de Octubre de 2014
	\end{center}



\newpage
\tableofcontents
\thispagestyle{empty}

\newpage
\renewcommand{\thepage}{\arabic{page}}
\setcounter{page}{1}

\section{INTRODUCCIÓN}
Un recinto clínico requiere generar un sistema de gestión de sus flujos de información. Se posee una base de datos existente, distribuida en diversos archivos de texto y se solicita a un estudiante de ingeniería informática generar un software, escrito en lenguaje de programación Scheme, bajo el paradigma de programación funcional, que sea capas de manejar un conjunto definido de funcionalidades. Para ello se le detallan reglas y formatos que deben seguir la compilación, ejecución y entrega de resultados de dicho software.
\\
\\
El presente documento resume el trabajo realizado en el desarrollo del software solicitado, a su vez explica cada funcionalidad implementada y los problemas que surgieron en el desarrollo de las mismas.  Este documento se compone de los siguientes capítulos:
\begin{itemize}
\item Descripción del problema, sección donde se especifican cada funcionalidad a desarrollar, reglas y formatos exigidos.
\item Descripción de solución, sección donde se detalla cada una de las funciones y/o procedimientos desarrollados para cumplir las operatividades exigidas.
\item Conclusiones, sección donde se realiza una visión global de lo aprendido, problemas y/o errores presentes en el desarrollo y el cumplimiento de los objetivos planteados.
\item Referencias, sección donde se informa cada una de las fuentes de información consultadas.
\item Anexo Codigo Fuente.
\end{itemize}

\section{DESCRIPCIÓN DEL PROBLEMA}

	\subsection{Motivación}
	Para el recinto clínico \textsl{FastClinic} es de imperiosa necesidad manejar su información de una manera sencilla, rápida y centralizada. Su antiguo sistema informático, sin soporte, ya no cubría sus necesidades y por es por esta razón que se ven en la necesidad de solicitar el desarrollo de un nuevo software, teniendo como base de registros archivos de textos heredados del antiguo sistema.

	\subsection{Especificaciones del proceso de atención}
	Para el desarrollo del nuevo sistema, la directiva de \textsl{FastClinic} nos ha entregado cierta información relevante, referente al proceso de atención de los pacientes y las relaciones existentes entre los distintos datos que se registran.
	\begin{enumerate}
	\item Un paciente es atendido por uno o más médicos y un médico puede atender múltiples pacientes.
	\item Un paciente puede llegar con uno o varios motivos de consultas.
	\item Un paciente puede tener uno o varios diagnósticos.
	\item Los diagnósticos son efectuados por un médico responsable en una fecha (dd/mm/aaaa) particular.
	\item Los pacientes pueden recibir distintos tratamientos recomendados por un doctor. Los tratamientos están normalizados para cada diagnóstico.
	\item Los diagnósticos pueden ser tratados con uno o más tratamientos. Los tratamientos pueden ser comunes para uno o más diagnósticos.
	\item Los pacientes reciben el alta de un diagnóstico en una fecha (dd/mm/aaaa) particular y bajo la responsabilidad de un médico (que puede ser diferente del médico que hizo el diagnóstico original).
	\end{enumerate}

	\subsection{Funcionalidades solicitadas}
	Dentro del enunciado del proyecto a desarrollar, se detallan diversas funcionalidades que pueden ser implementadas. No se establece que deben implementarse todas, quedando a criterio del tiempo del estudiante decidir cual programar. \\
	En el siguiente capitulo, se especifican las funcionalidades implementadas como parte de la solución. En el caso que se quiera conocer el abanico de funciones que podrían haber sido implementadas, favor dirigirse al enunciado del proyecto.

	\subsection{Dominio del problema}
	Por parte del estudiante de ingeniería, existe un control medio del problema ya que gracias a lo aprendido en cursos anteriores puede aplicar la lógica de resolución de problemas, identificando las similitudes fuertes o debiles entre las diversas funcionalidades solicitadas para, de esta forma, aprovechar de mejor manera el tiempo invertido. Por otra parte,  el manejo del lenguaje scheme es una falencia, ya que estos contenidos, aunque han sido entendidos en las clases de la asignatura, no se han interiorizado en gran medida. Notemos que el cambio de paradigma es algo complicado en primera instancia y esto, genera retrasos de implementación en el proyecto. 

%Nueva sección
\section{DESCRIPCIÓN DE LA SOLUCIÓN}
	\subsection{Propósito de la solución}
	El objetivo principal es generar un software de gestión que cumpla los estándares especificados por la gente de FastClinic y que sea capaz de entregar satisfacción al usuario, optimizando el tiempo invertido en los procesos productivos de dicha entidad.
	
	\subsection{Características de la solución}
	El software generado estará codificado en lenguaje Scheme, bajo las reglas de la programación funcional. Se debe destacar que dicho paradigma de programación utliza el tipo de dato abstracto Lista y operaciones recursivas como sus piedras angulares de implementación y procesamiento.\\
   
   Un tipo de dato abstracto, es una implementación de codigo que permite emular el comportamiento que tiene un "entidad" o "ser" a nivel real. En otras palabras, permite generar una versión programada de los atributos y las acciones que desempeña un determinado individuo. Para este caso, se han implementado los siguientes tipos de datos abstractos:
    \begin{enumerate}
	\item TDA Doctor
	\item TDA Paciente
	\item TDA Tratamiento
	\item TDA Diagnostico
	\item TDA DiagnosticoPaciente
	\item TDA TratamientoDiagnostico
	\item TDA TratamientosDiagnosticosPaciente
	\end{enumerate}

	\subsection{Desafios de la solución}
	El principal problema o desafio, para el estudiante de ingeniería informática, es poder interiorizar de buena forma los conceptos y la lógica de la programación funcional demostrando, dicho aprendizaje de contenidos, de forma practica a través de la creación del software solicitado.
	
	\subsection{Funcionalidades implementadas}
	A continuación se detallan cada una de las funcionalidades trabajadas. Se debe destacar que el total de funcionalidades implementadas son 9, distribuyendose de la siguiente forma:
	\begin{itemize}
	\item 5 Funcionalidades de nivel bajo.
	\begin{enumerate}
	\item obtenerNombrePaciente
	\item especialidad
	\item tratamientoRiesgoso
	\item cantidadMedicosTratamiento
	\item cantidadPacientesMedico
	\end{enumerate}
	\item 2 Funcionalidades de nivel medio.
	\begin{enumerate}
	\item tratamientoMasUsadoPorDiagnostico
	\item medicoMasAltas
	\end{enumerate}
	\item 2 Funcionalidades de nivel alto.
	\begin{enumerate}
	\item listarMedicosTratantesPaciente
	\item diagnosticoPacienteMedico
	\end{enumerate}
	\item Nota: Además de estas funcionalidades, se implementaron: 
	\begin{itemize}
	\item Selectores
	\item Constructores
	\item Modificadores
	\item Verificadores de pertenencias
	\end{itemize}
	\end{itemize}

Para ver en detalle las implementaciones desarrolladas, favor dirigirse al Anexo de codigo fuente, adjunto al final del rpesente documento.
	
\section{ANÁLISIS DE RESULTADOS}
En el desarrollo, solo se desarrollaron las funcionalidades que a juicio del estudiante fueran abordables con repsecto al tiempo reducido con el que disponia. Es por esta razón que los resultados, teniendo en cuenta unica y exclusivamente las funcionales trabajadas, no cumplen con lo esperado. Sin embargo, cada funcionalidad desarrollada, logra dar solución a lo que se solicitó por parte de FastClinic. \\

	\subsection{Alcance de la solución}
	En concreto, se alcanza a cumplir y sobrepasar en baja intensidad el mínimo exigido. Este alcance no satisface al estudiante, sin embargo, es realmente lo que se implementa.

\newpage
\section{CONCLUSIÓN}
Durante el desarrollo del software, se presentaron diversas complicaciones. Por una parte, los errores de redacción del enunciado del proyecto propiciaron confusiones del estudiante en cuestión y generaron retardos en la implementación de las funcionalidades solicitadas. Por otro lado, el cambio de paradigma es dificil de comprender, recordemos que el efecto paradigma se hace presente cada vez que intentamos mirar las cosas desde otra perspectiva y nos impide salir de las reglas actuales, en sintesis, nos cuesta reconocer otras formas de hacer las cosas. \\

En cuanto al lenguaje de programación Scheme, pese a ser un lenguaje antiguo no se queda atrás en poder de procesamiento. Esto puede verse muy bien en la aplicación que tuvo este lenguaje en manejo de grandes tamaños de datos. Su papel en la actualidad es un tanto desconocido, pero no hay duda de que es un lenguaje poderoso aunque de dificil lectura y comprensión. \\

Finalmente, se puede visualizar bajo la propia experiencia las diferencias entre los distintos paradigmas de programación, teniendo ahora como comparar objetivamente, en un ambiente real las ventajas y desventajas que nos otorga cada paradigma en cuestión.

\section{REFERENCIAS}
\begin{itemize}
\item Andrés Marzal, Isabel Garcia. (Año no informado). Introducción a la programación con Scheme. Departamento de lenguajes y sistemas informáticos: Universitat Jaume.
\item  Eduardo Navas (2010). Programando con Racket 5.
El Salvador: Universidad Centroamericana "Jose Simeón Cañas".
\end{itemize}

\newpage

\clearpage
\addcontentsline{toc}{section}{Referencias}
\bibliographystyle{unsrt}
\bibliography{references}

\end{document}
